\section{بازتعریف مسئله و مبانی نظری}

\begin{frame}{چرا این مسئله با ابزارهای استاندارد رگرسیون حل نمی‌شود؟}
  \begin{itemize}
    \item \textbf{نامتقارنی هزینه‌ی خطا:} کمبود گران‌تر از مازاد است — معیاری که خطای مثبت و
      منفی را یکسان می‌بیند (مثل \lr{MAE}، \lr{RMSE}) با فرض غلط $C_u=C_o$ کار می‌کند
    \item \textbf{خروجی مدل یک چندک است، نه یک عدد میانگین:} مدل باید $\hat\rho_\tau$ را بزند،
      نه $E[\rho]$ — پس معیار باید یک \lr{proper scoring rule} برای چندک باشد
    \item \textbf{ساختار پنل چندسری‌زمانی:} چند سلف × وعده × نوع غذا هم‌زمان و در طول زمان —
      نحوه‌ی اعتبارسنجی (زمانی، تفکیک بین سری‌ها) به‌اندازه‌ی فرمول معیار اهمیت دارد
    \item \textbf{متغیر هدف نرخی و کران‌دار:} $\rho\in[0,1]$ و در رکوردهای کم‌رزرو می‌تواند
      دقیقاً صفر شود — معیارهایی مثل \lr{MAPE} همین‌جا از کار می‌افتند
  \end{itemize}
  \note{این چهار ویژگی تعیین می‌کنند کدام معیارهای رایج قابل‌استفاده، کدام‌ها گمراه‌کننده و
  کدام‌ها فقط مکمل‌اند — جزئیات کامل در سند پشتیبان \lr{model-evaluation-metrics.md}. این اسلاید
  زمینه‌چینی برای مسئله‌ی روزنامه‌فروش در اسلاید بعد است؛ کوتاه رد شو، تفصیلش در بخش معیارهای
  ارزیابی می‌آید.}
\end{frame}

\begin{frame}{مسئله‌ی روزنامه‌فروش: تصمیم‌گیری تحت زیان نامتقارن}
  \begin{columns}
    \begin{column}{0.50\textwidth}
      \begin{itemize}
        \item \textbf{دوراهی تصمیم پخت (۲۰۰ رزرو):} ۲۰۰ پخت و ۱۵ غایب $\leftarrow$ هدررفت ۱۵ پرس ($5.6$ میلیون تومان)؛ اما ۱۸۵ پخت و ۱۹۰ متقاضی $\leftarrow$ کمبود ۵ پرس. کدام خطا پرهزینه‌تر است؟
        \item \textbf{چارچوب روزنامه‌فروش (\lr{Newsvendor}):} تصمیم‌گیری مقداری پیش از تقاضا تحت هزینه‌ی نامتقارن کمبود ($C_u$) در برابر مازاد ($C_o$).
        \item \textbf{چندک بحرانی بهینه:} پاسخ بهینه میانگین تقاضا نیست، چندک تراز است:
          \[
            \tau^{*} = \frac{C_u}{C_u + C_o}
          \]
      \end{itemize}
    \end{column}
    \begin{column}{0.48\textwidth}
      \begin{figure}
        \centering
        \includegraphics[width=\textwidth]{slides_03_asymmetric_cost}
      \end{figure}
    \end{column}
  \end{columns}
  \note{این همان مسئله‌ی کلاسیک Newsvendor در پژوهش در عملیات است.
  جمله‌ی کلیدی برای ارائه: اگر هزینه‌ی کمبود ۴ برابر مازاد باشد ($C_u/C_o = 4$)، مقدار بهینه‌ی
  پخت نقطه‌ای است که فقط در $20\%$ مواقع دچار کمبود شویم ($\tau^* = 0.20$).}
\end{frame}

\begin{frame}{شکست فرمول‌های کلاسیک: نوسان واقعی تا $15.6$ برابر فرضیات نظری است}
  \begin{itemize}
    \item \textbf{رد فرض استقلال افراد:} فرمول‌های تحلیلی غیبت را مستقل می‌دانند (توزیع دوجمله‌ای)، اما رفتار جمعی نوسان را تا $15.6$ برابر فراتر از تئوری می‌برد
    \item \textbf{ماندگاری نوسان پس از کنترل ساختاری:} حتی با کنترل سلف، وعده، تقویم و روز هفته، همچنان $7.6$ برابر نوسان اضافه باقی می‌ماند
    \item \textbf{پیامد راهبردی:} در میان ۷٬۵۷۹ سلول ناهمگن، چندک بهینه فرمول بسته ندارد و باید \textbf{مستقیماً از داده آموخته شود}
  \end{itemize}
  \begin{figure}
    \centering
    \includegraphics[width=0.55\textwidth]{4.8_variance_vs_res}
  \end{figure}
  \note{در رویکرد کلاسیک پژوهش عملیات، ابتدا یک توزیع پارامتری (مثل دوجمله‌ای با فرض استقلال
  افراد) برازش می‌شود و سپس چندک بهینه تحلیلی حساب می‌شود. اما داده‌ها نشان دادند فرض استقلال
  کاملاً باطل است و نوسان واقعی بین $3.7$ تا $15.6$ برابر فرمول نظری است (بیش‌پراکندگی صعودی).
  حتی کنترل کامل زمان و مکان هم این شکاف را پر نمی‌کند ($7.6$ برابر نوسان اضافه باقی می‌ماند).
  ضمناً ۷٬۵۷۹ سلول مستقل با شرایط متفاوت داریم. نتیجه: چندک نباید از فرمول بسته بیرون بیاید،
  باید مستقیماً از داده یاد گرفته شود. شاهد: \lr{F07}، \lr{F08}، \lr{F09}.}
\end{frame}

\begin{frame}{نیوزوندور داده‌محور و رگرسیون کوانتایل}
  \begin{itemize}
    \item به‌جای اینکه اول توزیع تقاضا را تخمین بزنیم و بعد عدد بهینه را حساب کنیم، مدل
      مستقیماً همان «عدد پخت» را از داده یاد می‌گیرد
    \item ابزار: رگرسیون کوانتایل — به خطای کم‌زدن وزن $\tau$ و به خطای زیادزدن وزن $1-\tau$
      می‌دهد
    \item این یعنی یادگیری مدل و حل مسئله‌ی روزنامه‌فروش، یک کار واحدند نه دو مرحله‌ی جدا
  \end{itemize}
  \begin{figure}
    \centering
    \includegraphics[width=0.68\textwidth]{slides_05_pinball}
  \end{figure}
  \note{مدل مستقیماً «عدد پخت» را می‌آموزد، نه «تقاضا» را. به همین دلیل مدل نهایی یک درخت
  گرادیان‌بوستینگ با هدف کوانتایلی است، همان انتخابی که در رقابت \lr{M5 Uncertainty} هم برنده
  شد. فرمول Pinball را بلند نگو، اسلاید پشتیبان دارد.}
\end{frame}
